% Final NN-G036 row exclusion and exact conclusion; root whole-proof PASS 2026-11-22T00:21:30Z.
\paragraph{The last mass-twelve row.}
The only row not covered by the preceding arguments is
\[
 w=12,\quad g_1=g_2=96,\quad s=768,\quad
 \delta_1=\delta_2=\psi,\quad u_1=u_2=10.
\]
Every effective $U_i\ge D$ of mass ten with the required
parity is $D+6$ or $D+2+2T_j$. In particular
$U_i(0)\in\{6,-2\}$, so both values are nonzero.

If $W(0)\ne0$, it is $8$ or $-8$ by the mass-twelve
classification. In the pair of cell identities
\[
 WG_1-U_1G_2=DW\psi,\qquad
 WG_2-U_2G_1=DW\psi,
\]
the coefficient determinant at zero is
$64-U_1(0)U_2(0)\ne0$, because the possible products
are $36,-12,4$. Hence both $G_i$ are divisible by $D$.
After cancelling this factor, the right sides are
$W\psi$, still divisible by $D$.
Applying the same invertible matrix at zero again
shows that both $G_i$ are divisible by $D^2$.
Since $WS=G_1G_2$ and $W(0)\ne0$, it follows that
$D^4\mid S$ integrally.
The quotient has even coefficients because $S\equiv0$
modulo two, but its value at four is
$768/4^4=3$. This is impossible.

It remains that $W=DH$, where $H(4)=3$,
$H\equiv1\pmod2$, and $H$ is either three
or primitive irreducible with leading coefficient two.
If $H$ is nonconstant, it cannot divide either
$U_i$, since three does not divide ten.
The common-factor argument of (G036.R2) gives
$H\mid G_1,G_2$ integrally.
If $H=3$, reduce both cell identities modulo three.
The reductions of the $U_i$ are nonzero, since their
masses are ten; the integral domain $\mathbb F_3[D]$
therefore forces the reductions of both $G_i$ to
vanish. Thus $H\mid G_1,G_2$ also in this case.
Write $G_i=HP_i$. Then
\[
 P_i(4)=32,\qquad
 P_1\equiv0,\quad P_2\equiv D^2(D-1)\pmod2,
\]
and the divided cell identities are
\[
 DH P_1-U_1P_2=D^2\psi,\qquad
 DH P_2-U_2P_1=D^2\psi.
\]
Because each $U_i(0)\ne0$, evaluation at zero
forces both $P_i(0)=0$.
Writing $P_i=DQ_i$ and cancelling one $D$ gives
\[
 DH Q_1-U_1Q_2=D\psi,\qquad
 DH Q_2-U_2Q_1=D\psi.
\]
Evaluation at zero again forces both $Q_i(0)=0$.
Consequently $P_i=D^2R_i$ with integral $R_i$.
Their masses are $R_i(4)=32/16=2$.
But $P_2\equiv D^2(D-1)$ forces
$R_2\equiv D-1\pmod2$, whose evaluation at four
is odd. This contradicts $R_2(4)=2$.
The last row is excluded for every grading support.
\hfill$\square$

\paragraph{Exact closed-table optimum and geometric conclusion.}
The complete necessary mass gate of NN-G035,
followed by all the uniform exclusions of NN-G036,
leaves exactly the unique $2744$-pair realization
of NN-G033--NN-G034. Therefore the minimum number
of cancelling trivial/sign pairs in the specified
symmetric closed-table presentation is exactly
$2744$, and its polynomial parameter realization
is unique within that presentation.
All cell degrees and contents, all grading support
locations, and all rational auxiliary parameters
were allowed throughout, subject only to integrality
and effectivity of the actual structure constants.
The final certificate projects every surviving
ordinary-dimension row; it does not enumerate
bounded grading degrees to justify the exclusions.

In the NN-G010 category this is an exact minimum
for actual effective graded sign-representation
objects composing by the fixed kernel table.
The grading is the diagonalizable-group dictionary
of SGA~3, I, 4.7.2--4.7.3, and composition is
the supplied archive cote~2, scan~17 transcription.
The theorem concerns this presentation, not all
possible effective multiplication tables.
The unrestricted extra-position optimization still
has the actual upper bound $2744$, the finite
equal-score cutoff of NN-G029/NN-G033, and the
total-mass FPT algorithm of NN-G032; its numerical
minimum remains open.
